\newpage
{\centering
\section{Appendice}
\addtocontents{toc}{\setcounter{tocdepth}{0}}
}
\medskip
{\color{gray}\hrule}
\medskip
\subsection{Calibrazione del magnete}
Sono qui riportati i dati sperimentali utilizzati per ricavare la curva di
calibrazione del magnete.
\subsubsection{Uniformità del campo}
\label{sec:unif}
Si riportano in seguito le misure del campo magnetico nello spazio tra le estensioni
polari a circa $I=2.3$A.
\begin{table}[ht!]
    \centering
    \begin{tabular}{|c|c|c|c|c|c|c|c|}
	\hline
	\textbf{B[mT]} & 145.85 & 150.43 & 168.49 & 142.42 & 149.49 & 155.11 & 138.31\\
	\hline
    \end{tabular} 
    \label{tab:B_unif}
\end{table}
Le misure sono state prese in cinque punti sul piano ortogonale al campo magnetico
a metà tra le estensioni polari e nei due punti centrali adiacenti ad esse.
\subsubsection{Misure di B al variare di I}
\label{sec:B_I}
Si riportano ora i dati utilizzati per ricavare la curva di calibrazione del magnete.
\begin{multicols}{2}
\begin{table}[H]
	\centering
    \begin{tabular}{|c|c|}
    \hline
    \textbf{I[A]} & \textbf{B[mT]} \\ \hline\hline
    $0.00\pm0.10$ &	$6  \pm15	$\\ \hline
    $1.15\pm0.13 $&	$82 \pm15 $		\\ \hline
    $1.95\pm0.15 $&	$135\pm15$	\\ \hline
    $3.04\pm0.18 $&	$197\pm15$	\\ \hline
    $3.9\pm0.2 	 $&	$259\pm15$	\\ \hline
    $5.1\pm0.2 	 $&	$334\pm17$	\\ \hline
    $6.4\pm0.3 	 $&	$420\pm20$	\\ \hline
    $7.0\pm0.3 	 $& 	$460\pm20$	\\ \hline
    $7.9\pm0.3   $&	$520\pm30$	\\ \hline
    $9.0\pm0.3   $&	$580\pm30$	\\ \hline
    $9.8\pm0.3   $&	$610\pm30$	\\ \hline
\end{tabular}
\caption{Dati B(I), salita n.1}
\end{table}
\begin{table}[H]
	\centering
    \begin{tabular}{|c|c|} \hline
\textbf{I[A]} & \textbf{B[mT]}	\\ \hline\hline
$9.8\pm0.3  $&	$610\pm	30$	\\ \hline
$9.0\pm0.3  $&	$580\pm	30$	\\ \hline
$8.1\pm0.3  $& $530\pm	30$	\\ \hline
$7.1\pm0.3  $& $470\pm	20$	\\ \hline
$6.0\pm0.3  $& $410\pm	20$	\\ \hline
$5.1\pm0.2  $& $341\pm	17$	\\ \hline
$4.0\pm0.2  $& $273\pm	15$	\\ \hline
$3.1\pm0.2  $& $210\pm	15$	\\ \hline
$1.95\pm0.18 $& $145\pm	15$	\\ \hline
$1.12\pm0.13$&	$88 \pm 15$	\\ \hline
$0.00\pm0.10  $& $8  \pm 15$	\\ \hline
\end{tabular}
\caption{Dati B(I), discesa n.1}
\end{table}
\end{multicols}
%\newpage 
\begin{multicols}{2}
\begin{table}[H]
	\centering
    \begin{tabular}{|c|c|}
    \hline
    \textbf{I[A]} & \textbf{B[mT]}	\\ \hline\hline
$0.00\pm0.10$&	$8  \pm	15$\\ \hline	
$0.96\pm0.12$&	$70 \pm 15$\\ \hline	
$2.30\pm0.16$&	$159\pm	15$\\ \hline	
$2.94\pm0.17$&	$192\pm	15$\\ \hline	
$3.9\pm	0.2 $&	$255\pm	15$\\ \hline	
$5.1\pm	0.2 $&  $335\pm	17$\\ \hline	
$6.2\pm	0.3 $&  $410\pm	20$\\ \hline	
$6.9\pm	0.3 $&   $460\pm20$\\ \hline	
$8.0\pm	0.3 $&  $520\pm	30$\\ \hline	
$9.0\pm	0.3 $&  $580\pm	30$\\ \hline	
$9.8\pm	0.3 $&  $610\pm	30$\\ \hline	
\end{tabular}
\caption{Dati B(I), salita n.2}
\end{table}
\begin{table}[H]
	\centering
    \begin{tabular}{|c|c|} \hline
\textbf{I[A]} & \textbf{B[mT]}	\\ \hline\hline
$9.8\pm0.3 $&	$610\pm30. $	\\ \hline
$9.0\pm0.3 $&  $580\pm30. $\\ \hline    
$8.1\pm0.3 $&  $540\pm30. $\\ \hline    
$6.9\pm0.3 $&    $460\pm20. $\\ \hline    
$6.0\pm0.3 $&    $400\pm20. $\\ \hline    
$5.1\pm0.2 $&    $342\pm17. $\\ \hline    
$4.0\pm0.2 $&    $274\pm15. $\\ \hline    
$3.13\pm0.18$&	$213\pm15. $  \\ \hline  
$2.03\pm0.15$&	$150\pm15.$	\\ \hline
$0.80\pm0.12$&	$65\pm15. $  \\ \hline  
$0.00\pm0.10$&	$8\pm15.  $  \\ \hline  
\end{tabular}
\caption{Dati B(I), discesa n.2}
\end{table}
\end{multicols}

\subsection{Misura dell'effetto Zeeman}
Sono qui riportati i dati sperimentali utilizzati per stimare il magnetone di Bohr.
\subsubsection{Geometria longitudinale}
\label{sec:longit}
Si riportano i valori di $\langle\delta\rangle$, $\langle\Delta\rangle$ e del loro
rapporto al variare dell'intensità di corrente $I$, con il suo corrispettivo valore
di campo magnetico $B$ per il setup in configurazione longitudinale

\begin{table}[H]
    \centering
    \begin{tabular}{|c|c|c|c|c|} \hline
	\textbf{I[A]} & \textbf{B[mT]} & $\mathbf{\langle\delta\rangle}\text{\textbf{[}}\mathbf{\mu}\text{\textbf{m}}^2\textbf{]}$  & $\mathbf{\langle\Delta\rangle}\text{\textbf{[}}\mathbf{\mu}\text{\textbf{m}}^2\textbf{]}$  & $\mathbf{\langle\delta\rangle/\langle\Delta\rangle}$ \\ \hline\hline
$0.0\pm0.1$ & $12\pm15$ & $0\pm100$ 	& $17500\pm100$ & $0\pm0.006$ \\ \hline
$2.54\pm0.16$ & $174\pm15$ &$2370\pm170$ & $17190\pm170$ &$0.138\pm0.009$ \\ \hline
$3.24\pm0.18$ & $219\pm15$ &$3140\pm70$ 	& $17200\pm90$ &$0.182\pm0.004$ \\ \hline
$4.0\pm0.2$ & $269\pm15$ &$4240\pm140$ 	& $17330\pm110$ &$0.245\pm0.008$ \\ \hline
$5.0\pm0.2$ & $334\pm15$ &$5020\pm120$ 	& $17300\pm110$ &$0.290\pm0.007$ \\ \hline
$6.1\pm0.3$ & $400\pm15$&$5890\pm130$ 	& $17660\pm180$ &$0.334\pm0.008$ \\ \hline
$7.2\pm0.3$ & $474\pm15$ &$6950\pm120$ 	& $17270\pm130$ &$0.402\pm0.007$ \\ \hline
$8.0\pm0.3$ & $520\pm15$ &$7880\pm70$ 	& $17310\pm60$ &$0.455\pm0.004$ \\ \hline
$9.0\pm0.3$ & $584\pm15$ &$8410\pm70$ 	& $17250\pm50$ &$0.488\pm0.004$ \\ \hline
\end{tabular}
\label{tab:longit}
\end{table}

\subsubsection{Geometria trasversale}
\label{sec:trasv}
Si riportano i valori di $\langle\delta\rangle$, $\langle\Delta\rangle$ e del loro
rapporto al variare dell'intensità di corrente $I$, con il suo corrispettivo valore
di campo magnetico $B$ per il setup in configurazione trasversale
\begin{table}[H]
    \centering
    \begin{tabular}{|c|c|c|c|c|} \hline

\textbf{I[A]} & \textbf{B[mT]} & $\mathbf{\langle\delta\rangle}\text{\textbf{[}}\mathbf{\mu}\text{\textbf{m}}^2\textbf{]}$  & $\mathbf{\langle\Delta\rangle}\text{\textbf{[}}\mathbf{\mu}\text{\textbf{m}}^2\textbf{]}$  & $\mathbf{\langle\delta\rangle/\langle\Delta\rangle}$ \\ \hline\hline
$0.0	\pm0.1$&$12	\pm8	$&$0	\pm90	$&$17300\pm100	$&$0.000	\pm0.005$  \\ \hline
$5.7	\pm0.2$&$373	\pm15	$&$2650	\pm50	$&$17380\pm60	$&$0.152	\pm0.003$  \\ \hline
$6.3	\pm0.3$&$414	\pm15	$&$3170	\pm70	$&$17220\pm60	$&$0.184	\pm0.004$  \\ \hline
$7.0	\pm0.3$&$462	\pm15	$&$3320	\pm60	$&$17160\pm50	$&$0.194	\pm0.003$  \\ \hline
$7.6	\pm0.3$&$497	\pm15	$&$3710	\pm30	$&$17510\pm20	$&$0.2118	\pm0.0017$ \\ \hline
$9.0	\pm0.3$&$588	\pm15	$&$4200	\pm100	$&$17290\pm130	$&$0.242	\pm0.006$  \\ \hline
$9.6	\pm0.3$&$624	\pm16	$&$4300	\pm90	$&$17600\pm80	$&$0.244	\pm0.005$  \\ \hline
\end{tabular}
    \label{tab:trasv_up}
    \caption{Dati per la transizione con $\Delta m_l = +1$}
\end{table}
\begin{table}[H]
    \centering
    \begin{tabular}{|c|c|c|c|c|} \hline

\textbf{I[A]} & \textbf{B[mT]} & $\mathbf{\langle\delta\rangle}\text{\textbf{[}}\mathbf{\mu}\text{\textbf{m}}^2\textbf{]}$  & $\mathbf{\langle\Delta\rangle}\text{\textbf{[}}\mathbf{\mu}\text{\textbf{m}}^2\textbf{]}$  & $\mathbf{\langle\delta\rangle/\langle\Delta\rangle}$ \\ \hline\hline
$0.0	\pm0.1	$&$12\pm15   $&$0	\pm90	$&$17300\pm100	$&$0.000	\pm0.005$\\ \hline
$5.7	\pm0.2	$&$373\pm15 $&$2600	\pm100	$&$17090\pm90	$&$-0.154	\pm0.006$\\ \hline
$6.3	\pm0.3	$&$414\pm15 $&$2840	\pm70	$&$17100\pm60	$&$-0.166	\pm0.004$\\ \hline
$7.0	\pm0.3	$&$462\pm15 $&$3360	\pm60	$&$17080\pm50	$&$-0.197	\pm0.003$\\ \hline
$7.6	\pm0.3	$&$497\pm15 $&$3670	\pm40	$&$17540\pm30	$&$-0.209	\pm0.002$\\ \hline
$9.0	\pm0.3	$&$588\pm15 $&$3980	\pm50	$&$17190\pm90	$&$-0.231	\pm0.003$\\ \hline
$9.6	\pm0.3	$&$624\pm16 $&$4410	\pm130	$&$17310\pm110	$&$-0.255	\pm0.008$\\ \hline
\end{tabular}
    \label{tab:trasv_down}
    \caption{Dati per la transizione con $\Delta m_l = -1$}
\end{table}
